\chapter{Graphical Solution and the Geometry of Linear Programs}

A linear program in two variables can be solved completely with pencil,
paper, and a basic calculator. The constraints cut a convex polygon out of
the first quadrant, and the optimum of a linear objective (if one exists) is always attained at a corner of that polygon. The graphical method therefore reduces to algebra: intersect the
constraint lines pairwise, keep those intersection points that satisfy every
constraint, and evaluate the objective at each of them. This chapter drills
that routine together with the geometric vocabulary (feasible region,
corner point, iso-profit line, binding constraint) and the two ways an LP
can fail to have an optimal corner: unboundedness and infeasibility. The
same geometry, in $n$ dimensions, is exactly what makes the simplex method
of the later chapters work.

\section{Core Rules}

\begin{rulebox}{The feasible region is an intersection of half-planes}
A linear constraint $a_{1}x_{1}+a_{2}x_{2}\le b$ describes a
\emph{half-plane}: the boundary line
$a_{1}x_{1}+a_{2}x_{2}=b$ together with all points on one side of it. To
decide which side, test a single point not on the line (the origin, whenever
it is not on the line): if the test point satisfies the inequality, its side
is the feasible one. The \emph{feasible region}
of an LP is the intersection of all its half-planes, including the sign
constraints $x_{1}\ge 0$ and $x_{2}\ge 0$. An intersection of half-planes is
always convex; in two variables it is a convex polygon, possibly
unbounded, possibly empty.
\end{rulebox}

\begin{rulebox}{Corner points and how to compute them}
A \emph{corner point} (vertex) of the feasible region is a
feasible point at which two constraint boundary lines with different
directions meet. With $m$ structural constraints plus the two sign
constraints there are at most $\binom{m+2}{2}$ candidate points. Procedure:
\begin{enumerate}[leftmargin=1.8em,itemsep=0.3ex]
  \item Take every pair of boundary lines (structural constraints as
        equalities, plus $x_{1}=0$ and $x_{2}=0$).
  \item Solve each $2\times 2$ linear system; parallel lines give no
        candidate.
  \item Check the candidate against \emph{all remaining} constraints,
        including non-negativity. Discard it if any constraint is violated.
\end{enumerate}
The surviving candidates are precisely the corners of the feasible region.
\end{rulebox}

\begin{rulebox}{Iso-profit lines and the corner theorem}
The level sets of a linear objective $z=c_{1}x_{1}+c_{2}x_{2}$ are the
parallel lines $c_{1}x_{1}+c_{2}x_{2}=z$, called \emph{iso-profit lines} when maximizing and \emph{iso-cost lines} when minimizing. Shifting the line in the direction
of the vector $(c_{1},c_{2})$ increases $z$; maximizing means sliding the
line as far as possible in that direction while it still touches the
feasible region. \textbf{Corner theorem}: if the LP has an
optimal solution at all, then at least one corner is optimal; a bounded,
non-empty feasible region guarantees this. If \emph{two} corners are both
optimal, the objective is constant along the edge joining them, and every
point of that edge is optimal as well.
\end{rulebox}

\begin{rulebox}{Recipe: solving a two-variable LP by corner enumeration}
\begin{enumerate}[leftmargin=1.8em,itemsep=0.3ex]
  \item List all boundary lines and compute all pairwise intersection
        points.
  \item Discard every candidate that violates some constraint; the rest are
        the corners.
  \item Evaluate the objective $z$ at every corner.
  \item Maximization: the largest value wins; minimization: the smallest.
  \item If two corners tie for the optimal value, report the whole edge
        between them as optimal. If the region is unbounded, first check
        whether some feasible direction improves $z$ without limit.
\end{enumerate}
\end{rulebox}

\begin{rulebox}{Binding and slack constraints}
At a feasible point $(x_{1},x_{2})$ the constraint
$a_{1}x_{1}+a_{2}x_{2}\le b$ has \emph{slack}
\[
  s \;=\; b-a_{1}x_{1}-a_{2}x_{2}\;\ge\;0 .
\]
The constraint is \emph{binding} if $s=0$, i.e.\ the point
lies on the boundary line, and \emph{slack} if $s>0$.
At every corner of a two-variable LP at least two constraints are binding.
Economically, the constraints binding at the optimum are the scarce
resources; a slack constraint has leftover capacity equal to its slack
value.
\end{rulebox}

\begin{rulebox}{Special cases: unbounded and infeasible LPs}
\begin{itemize}[leftmargin=1.6em,itemsep=0.3ex]
  \item \textbf{Unbounded LP}: the feasible region
        contains a ray along which the objective improves without limit,
        no finite optimum exists. To certify it, exhibit feasible points
        $(x_{1}(t),x_{2}(t))$ for all $t\ge 0$ with $z\to\infty$ (or
        $-\infty$ when minimizing).
  \item \textbf{Infeasible LP}: the constraints
        contradict each other and the feasible region is empty,
        typically two parallel half-planes facing away from each other.
  \item An unbounded \emph{region} does not by itself make the LP
        unbounded: minimizing a positive-coefficient objective over a
        region that opens up and to the right still has a corner minimum,
        because moving into the region only increases the objective.
\end{itemize}
\end{rulebox}

\section{Practice Problems}

Work every problem by corner enumeration: list the candidate intersection
points, discard the infeasible ones, evaluate the objective at every
corner, and state the optimal value, the optimal point(s), and the
constraints binding there. All variables are non-negative:
$x_{1},x_{2}\ge 0$ holds in every problem, in addition to the constraints
listed.

\begin{problems}
  \item Solve $\max\ z=3x_{1}+2x_{2}$ subject to $x_{1}+x_{2}\le 4$.
  \item List all corner points of the region
        $\{(x_{1},x_{2}):x_{1}+x_{2}\le 5,\ 2x_{1}+x_{2}\le 8,\
        x_{1},x_{2}\ge 0\}$. For every pair of boundary lines state whether
        the intersection point is feasible.
  \item Maximize $z=5x_{1}+4x_{2}$ over the region of Problem~2.
  \item Solve $\max\ z=x_{1}+x_{2}$ subject to $2x_{1}+x_{2}\le 6$ and
        $x_{1}+2x_{2}\le 6$.
  \item Solve $\max\ z=4x_{1}+3x_{2}$ subject to $2x_{1}+x_{2}\le 10$ and
        $x_{1}+3x_{2}\le 15$.
  \item Solve $\max\ z=x_{1}+2x_{2}$ subject to $x_{1}+x_{2}\le 6$,
        $x_{1}\le 4$, $x_{2}\le 5$. At the optimum, classify every
        constraint as binding or slack and give the slack values.
  \item Solve $\max\ z=2x_{1}+x_{2}$ subject to $x_{1}+x_{2}\le 4$ and
        $2x_{1}+x_{2}\le 6$, and describe the \emph{complete} set of
        optimal solutions.
  \item Decide whether $\max\ z=x_{1}+x_{2}$ subject to $-x_{1}+x_{2}\le 1$
        has a finite optimum. If not, exhibit a feasible ray along which
        $z\to\infty$.
  \item Show that the constraint set $x_{1}+x_{2}\le 2$,
        $x_{1}+x_{2}\ge 5$, $x_{1},x_{2}\ge 0$ is infeasible, and describe
        the geometry of the two structural half-planes.
  \item Solve $\min\ z=2x_{1}+3x_{2}$ subject to $x_{1}+x_{2}\ge 4$.
        Explain why a corner minimum exists even though the feasible region
        is unbounded.
  \item Find \emph{all} optimal solutions of $\max\ z=x_{1}+2x_{2}$
        subject to $x_{1}+2x_{2}\le 10$.
  \item A workshop builds tables ($x_{1}$) and chairs ($x_{2}$). A table
        uses $2$ units of wood and $1$ hour of labour, a chair uses $1$
        unit of wood and $1$ hour of labour; $10$ units of wood and $6$
        hours of labour are available. Profits are $30$ per table and $20$
        per chair. Find the profit-maximizing production plan and state
        which resources are fully used.
  \item Find all corner points of the region
        $R=\{(x_{1},x_{2}):x_{1}+2x_{2}\le 8,\ 3x_{1}+2x_{2}\le 12,\
        x_{1},x_{2}\ge 0\}$, listing the infeasible candidate
        intersections as well.
  \item Maximize $z=3x_{1}+5x_{2}$ over the region $R$ of Problem~13.
  \item Maximize $z=3x_{1}+2x_{2}$ over the region $R$ of Problem~13. Is
        the optimal solution unique?
  \item A farmer has $40$ hectares of land and $60$ tons of fertilizer.
        One hectare of wheat ($x_{1}$) needs $2$ tons of fertilizer and
        yields profit $5$; one hectare of barley ($x_{2}$) needs $1$ ton
        and yields profit $4$. Maximize total profit.
  \item Solve $\max\ z=x_{1}$ subject to $x_{1}+x_{2}\le 4$ and
        $x_{1}-x_{2}\le 2$: that is, find the easternmost point of the
        feasible region.
  \item For $\max\ z=cx_{1}+x_{2}$ subject to $x_{1}+x_{2}\le 1$ with a
        parameter $c\ge 0$, find all values of $c$ for which $(1,0)$ is
        optimal. For which $c$ is it the \emph{unique} optimum?
\end{problems}

\section{Solutions}

\begin{solutions}
  \item Boundary lines: $x_{1}=0$, $x_{2}=0$, and (C1) $x_{1}+x_{2}=4$.
        Pairwise intersections: $(0,0)$; $x_{2}=0$ with (C1) gives
        $(4,0)$; $x_{1}=0$ with (C1) gives $(0,4)$. All three satisfy
        every constraint, so the corners are $(0,0)$, $(4,0)$, $(0,4)$.
        Objective values:
        \[
          z(0,0)=0,\qquad z(4,0)=12,\qquad z(0,4)=8 .
        \]
        Optimum $z^{*}=12$ at $(4,0)$. Binding there: (C1) (since
        $4+0=4$) and $x_{2}\ge 0$; the constraint $x_{1}\ge 0$ is slack.

  \item Label (C1) $x_{1}+x_{2}\le 5$ and (C2) $2x_{1}+x_{2}\le 8$. The
        six candidate intersections:
        \begin{itemize}[leftmargin=1.6em,itemsep=0.3ex]
          \item $x_{1}=0$, $x_{2}=0$: $(0,0)$; satisfies (C1), (C2).
                Corner.
          \item $x_{1}=0$ with (C1): $(0,5)$; check (C2):
                $2(0)+5=5\le 8$. Corner.
          \item $x_{1}=0$ with (C2): $(0,8)$; check (C1): $0+8=8>5$.
                Infeasible.
          \item $x_{2}=0$ with (C1): $(5,0)$; check (C2):
                $2(5)+0=10>8$. Infeasible.
          \item $x_{2}=0$ with (C2): $(4,0)$; check (C1): $4+0=4\le 5$.
                Corner.
          \item (C1) with (C2): subtracting $x_{1}+x_{2}=5$ from
                $2x_{1}+x_{2}=8$ gives $x_{1}=3$, hence $x_{2}=2$; both
                coordinates non-negative. Corner.
        \end{itemize}
        Corners: $(0,0)$, $(4,0)$, $(3,2)$, $(0,5)$.

  \item Using the corner list from Problem~2 and $z=5x_{1}+4x_{2}$:
        \[
          z(0,0)=0,\quad z(4,0)=20,\quad z(3,2)=15+8=23,\quad
          z(0,5)=20 .
        \]
        Optimum $z^{*}=23$ at $(3,2)$. Binding there: both structural
        constraints, since $3+2=5$ and $2(3)+2=8$; both sign constraints
        are slack.

  \item Label (C1) $2x_{1}+x_{2}=6$ and (C2) $x_{1}+2x_{2}=6$.
        Candidates: $(0,0)$ feasible; $x_{2}=0$ with (C1): $(3,0)$,
        check (C2): $3\le 6$, corner; $x_{2}=0$ with (C2): $(6,0)$,
        check (C1): $12>6$, infeasible; $x_{1}=0$ with (C1): $(0,6)$,
        check (C2): $12>6$, infeasible; $x_{1}=0$ with (C2): $(0,3)$,
        check (C1): $3\le 6$, corner. (C1) with (C2): doubling (C2)
        gives $2x_{1}+4x_{2}=12$; subtracting (C1) leaves $3x_{2}=6$, so
        $x_{2}=2$ and $x_{1}=2$: corner $(2,2)$. Evaluate
        $z=x_{1}+x_{2}$:
        \[
          z(0,0)=0,\quad z(3,0)=3,\quad z(2,2)=4,\quad z(0,3)=3 .
        \]
        Optimum $z^{*}=4$ at $(2,2)$; both structural constraints are
        binding ($2\cdot 2+2=6$ and $2+2\cdot 2=6$).

  \item Label (C1) $2x_{1}+x_{2}=10$ and (C2) $x_{1}+3x_{2}=15$.
        Candidates: $(0,0)$ feasible; $x_{2}=0$ with (C1): $(5,0)$,
        check (C2): $5\le 15$, corner; $x_{2}=0$ with (C2): $(15,0)$,
        check (C1): $30>10$, infeasible; $x_{1}=0$ with (C1): $(0,10)$,
        check (C2): $30>15$, infeasible; $x_{1}=0$ with (C2): $(0,5)$,
        check (C1): $5\le 10$, corner. (C1) with (C2): from (C1),
        $x_{2}=10-2x_{1}$; substituting into (C2),
        \[
          x_{1}+3(10-2x_{1})=15
          \;\Longrightarrow\; -5x_{1}=-15
          \;\Longrightarrow\; x_{1}=3,\ x_{2}=4 .
        \]
        Corners: $(0,0)$, $(5,0)$, $(3,4)$, $(0,5)$. Evaluate
        $z=4x_{1}+3x_{2}$:
        \[
          z(0,0)=0,\quad z(5,0)=20,\quad z(3,4)=12+12=24,\quad
          z(0,5)=15 .
        \]
        Optimum $z^{*}=24$ at $(3,4)$; binding: (C1) ($6+4=10$) and
        (C2) ($3+12=15$).

  \item Label (C1) $x_{1}+x_{2}\le 6$, (C2) $x_{1}\le 4$, (C3)
        $x_{2}\le 5$. Candidates (pairs of parallel or identical
        direction, such as $x_{1}=0$ with (C2), give no intersection):
        $(0,0)$ feasible; (C2) with $x_{2}=0$: $(4,0)$, check (C1):
        $4\le 6$, corner; (C1) with $x_{2}=0$: $(6,0)$, violates (C2),
        infeasible; (C3) with $x_{1}=0$: $(0,5)$, check (C1): $5\le 6$,
        corner; (C1) with $x_{1}=0$: $(0,6)$, violates (C3), infeasible;
        (C1) with (C2): $(4,2)$, check (C3): $2\le 5$, corner; (C1) with
        (C3): $(1,5)$, check (C2): $1\le 4$, corner; (C2) with (C3):
        $(4,5)$, check (C1): $9>6$, infeasible. Corners: $(0,0)$,
        $(4,0)$, $(4,2)$, $(1,5)$, $(0,5)$. Evaluate $z=x_{1}+2x_{2}$:
        \[
          z(0,0)=0,\quad z(4,0)=4,\quad z(4,2)=8,\quad z(1,5)=11,\quad
          z(0,5)=10 .
        \]
        Optimum $z^{*}=11$ at $(1,5)$. Classification there: (C1)
        binding ($1+5=6$, slack $0$); (C3) binding ($x_{2}=5$, slack
        $0$); (C2) slack with slack value $4-1=3$; $x_{1}\ge 0$ slack
        (value $1$), $x_{2}\ge 0$ slack (value $5$).

  \item Label (C1) $x_{1}+x_{2}\le 4$, (C2) $2x_{1}+x_{2}\le 6$.
        Candidates: $(0,0)$ feasible; $x_{2}=0$ with (C1): $(4,0)$,
        check (C2): $8>6$, infeasible; $x_{2}=0$ with (C2): $(3,0)$,
        check (C1): $3\le 4$, corner; $x_{1}=0$ with (C1): $(0,4)$,
        check (C2): $4\le 6$, corner; $x_{1}=0$ with (C2): $(0,6)$,
        check (C1): $6>4$, infeasible; (C1) with (C2): subtracting
        gives $x_{1}=2$, $x_{2}=2$, corner. Corners: $(0,0)$, $(3,0)$,
        $(2,2)$, $(0,4)$. Evaluate $z=2x_{1}+x_{2}$:
        \[
          z(0,0)=0,\quad z(3,0)=6,\quad z(2,2)=6,\quad z(0,4)=4 .
        \]
        The maximum $z^{*}=6$ is attained at \emph{both} $(3,0)$ and
        $(2,2)$: the iso-profit lines $2x_{1}+x_{2}=z$ are parallel to
        the binding constraint line $2x_{1}+x_{2}=6$. By the corner
        theorem the whole edge between the two corners is optimal:
        \[
          \{(3-t,\,2t):t\in[0,1]\},
        \]
        every point of which satisfies $2x_{1}+x_{2}=6$.

  \item The region $\{-x_{1}+x_{2}\le 1,\ x_{1},x_{2}\ge 0\}$ has
        corners $(0,0)$ and $(0,1)$ only (the pair $x_{2}=0$ with
        $-x_{1}+x_{2}=1$ gives $x_{1}=-1$, infeasible), but it is
        unbounded: for every $t\ge 0$ the point $(t,0)$ satisfies
        $-t+0\le 1$ and the sign constraints. Along this ray
        $z=x_{1}+x_{2}=t\to\infty$, so the LP is \emph{unbounded}: no
        finite optimum exists. This shows why corner enumeration alone
        is not conclusive over an unbounded region: one must also
        check improving directions.

  \item Suppose $(x_{1},x_{2})$ were feasible. Then simultaneously
        $x_{1}+x_{2}\le 2$ and $x_{1}+x_{2}\ge 5$, i.e.\
        $5\le x_{1}+x_{2}\le 2$, a contradiction. The region is
        empty and the LP infeasible. Geometrically the two structural
        constraints demand the half-plane on or below the line
        $x_{1}+x_{2}=2$ and the half-plane on or above the parallel
        line $x_{1}+x_{2}=5$; the two half-planes face away from each
        other across a gap and share no point.

  \item The region $\{x_{1}+x_{2}\ge 4,\ x_{1},x_{2}\ge 0\}$ is
        unbounded (it contains $(t,t)$ for all $t\ge 2$). Its corners:
        $x_{2}=0$ with $x_{1}+x_{2}=4$ gives $(4,0)$; $x_{1}=0$ with
        the same line gives $(0,4)$; the origin violates
        $x_{1}+x_{2}\ge 4$ and is infeasible. Both objective
        coefficients of $z=2x_{1}+3x_{2}$ are positive, so moving
        deeper into the region (increasing either variable) can only
        increase $z$; the minimum therefore lies on the boundary and,
        by the corner theorem, at a corner. Evaluate:
        \[
          z(4,0)=8,\qquad z(0,4)=12 .
        \]
        Minimum $z^{*}=8$ at $(4,0)$; binding there:
        $x_{1}+x_{2}\ge 4$ and $x_{2}\ge 0$.

  \item Corners of $\{x_{1}+2x_{2}\le 10,\ x_{1},x_{2}\ge 0\}$:
        $(0,0)$, $(10,0)$, $(0,5)$. Evaluate $z=x_{1}+2x_{2}$:
        \[
          z(0,0)=0,\qquad z(10,0)=10,\qquad z(0,5)=10 .
        \]
        The two corners $(10,0)$ and $(0,5)$ tie at $z^{*}=10$: the
        objective is the left-hand side of the constraint, so the
        iso-profit line for $z=10$ \emph{coincides} with the constraint
        line. The optimal set is the entire edge
        \[
          \{(x_{1},x_{2}):x_{1}+2x_{2}=10,\ x_{1},x_{2}\ge 0\}
          =\{(10-2t,\,t):t\in[0,5]\} .
        \]

  \item Model: $\max\ z=30x_{1}+20x_{2}$ subject to wood (C1)
        $2x_{1}+x_{2}\le 10$, labour (C2) $x_{1}+x_{2}\le 6$,
        $x_{1},x_{2}\ge 0$. Candidates: $(0,0)$ feasible; $x_{2}=0$
        with (C1): $(5,0)$, check (C2): $5\le 6$, corner; $x_{2}=0$
        with (C2): $(6,0)$, check (C1): $12>10$, infeasible; $x_{1}=0$
        with (C1): $(0,10)$, check (C2): $10>6$, infeasible; $x_{1}=0$
        with (C2): $(0,6)$, check (C1): $6\le 10$, corner; (C1) with
        (C2): subtracting gives $x_{1}=4$, hence $x_{2}=2$, corner.
        Corners: $(0,0)$, $(5,0)$, $(4,2)$, $(0,6)$. Profits:
        \[
          z(0,0)=0,\quad z(5,0)=150,\quad z(4,2)=120+40=160,\quad
          z(0,6)=120 .
        \]
        Optimal plan: $4$ tables and $2$ chairs with profit
        $z^{*}=160$. Both resources are fully used (binding): wood
        $2(4)+2=10$ and labour $4+2=6$.

  \item Label (C1) $x_{1}+2x_{2}=8$ and (C2) $3x_{1}+2x_{2}=12$.
        Candidates: $(0,0)$ feasible, corner; $x_{2}=0$ with (C1):
        $(8,0)$, check (C2): $24>12$, infeasible; $x_{2}=0$ with (C2):
        $(4,0)$, check (C1): $4\le 8$, corner; $x_{1}=0$ with (C1):
        $(0,4)$, check (C2): $8\le 12$, corner; $x_{1}=0$ with (C2):
        $(0,6)$, check (C1): $12>8$, infeasible; (C1) with (C2):
        subtracting (C1) from (C2) gives $2x_{1}=4$, so $x_{1}=2$ and,
        from (C1), $2x_{2}=6$, $x_{2}=3$: corner $(2,3)$. Corners of
        $R$: $(0,0)$, $(4,0)$, $(2,3)$, $(0,4)$.

  \item Evaluate $z=3x_{1}+5x_{2}$ at the corners of $R$:
        \[
          z(0,0)=0,\quad z(4,0)=12,\quad z(2,3)=6+15=21,\quad
          z(0,4)=20 .
        \]
        Optimum $z^{*}=21$ at $(2,3)$; binding: both structural
        constraints ($2+6=8$ and $6+6=12$).

  \item Evaluate $z=3x_{1}+2x_{2}$ at the corners of $R$:
        \[
          z(0,0)=0,\quad z(4,0)=12,\quad z(2,3)=6+6=12,\quad
          z(0,4)=8 .
        \]
        The maximum $z^{*}=12$ is attained at both $(4,0)$ and
        $(2,3)$, so the optimum is \emph{not} unique: the iso-profit
        lines $3x_{1}+2x_{2}=z$ are parallel to the constraint
        $3x_{1}+2x_{2}=12$, which is binding along the whole optimal
        edge
        \[
          \{(4-2t,\,3t):t\in[0,1]\} .
        \]

  \item Model: $\max\ z=5x_{1}+4x_{2}$ subject to land (C1)
        $x_{1}+x_{2}\le 40$, fertilizer (C2) $2x_{1}+x_{2}\le 60$,
        $x_{1},x_{2}\ge 0$. Candidates: $(0,0)$ feasible; $x_{2}=0$
        with (C1): $(40,0)$, check (C2): $80>60$, infeasible;
        $x_{2}=0$ with (C2): $(30,0)$, check (C1): $30\le 40$, corner;
        $x_{1}=0$ with (C1): $(0,40)$, check (C2): $40\le 60$, corner;
        $x_{1}=0$ with (C2): $(0,60)$, check (C1): $60>40$,
        infeasible; (C1) with (C2): subtracting gives $x_{1}=20$,
        hence $x_{2}=20$, corner. Corners: $(0,0)$, $(30,0)$,
        $(20,20)$, $(0,40)$. Profits:
        \[
          z(0,0)=0,\quad z(30,0)=150,\quad z(20,20)=100+80=180,\quad
          z(0,40)=160 .
        \]
        Optimum: $20$ hectares of wheat and $20$ of barley, profit
        $z^{*}=180$. Both land ($20+20=40$) and fertilizer
        ($40+20=60$) are binding.

  \item Label (C1) $x_{1}+x_{2}\le 4$, (C2) $x_{1}-x_{2}\le 2$.
        Candidates: $(0,0)$ feasible; $x_{2}=0$ with (C1): $(4,0)$,
        check (C2): $4>2$, infeasible; $x_{2}=0$ with (C2): $(2,0)$,
        check (C1): $2\le 4$, corner; $x_{1}=0$ with (C1): $(0,4)$,
        check (C2): $-4\le 2$, corner; $x_{1}=0$ with (C2): $(0,-2)$,
        violates $x_{2}\ge 0$, infeasible; (C1) with (C2): adding the
        two equations gives $2x_{1}=6$, so $x_{1}=3$, $x_{2}=1$,
        corner. Corners: $(0,0)$, $(2,0)$, $(3,1)$, $(0,4)$. Evaluate
        $z=x_{1}$:
        \[
          z(0,0)=0,\quad z(2,0)=2,\quad z(3,1)=3,\quad z(0,4)=0 .
        \]
        Optimum $z^{*}=3$ at $(3,1)$: the easternmost corner, where
        (C1) and (C2) are both binding ($3+1=4$, $3-1=2$).

  \item Corners of $\{x_{1}+x_{2}\le 1,\ x_{1},x_{2}\ge 0\}$: $(0,0)$,
        $(1,0)$, $(0,1)$. Objective values: $z(0,0)=0$, $z(1,0)=c$,
        $z(0,1)=1$. The corner $(1,0)$ is optimal exactly when its
        value beats the other corners: $c\ge 1$ (which also gives
        $c\ge 0$). For $c>1$ the optimum is unique; for $c=1$ the
        corners $(1,0)$ and $(0,1)$ tie, the iso-profit lines
        $x_{1}+x_{2}=z$ are parallel to the constraint, and the whole
        edge $x_{1}+x_{2}=1$, $x_{1},x_{2}\ge 0$ is optimal.
\end{solutions}

\section{Additional Problems}

As before, $x_{1},x_{2}\ge 0$ holds in every problem in addition to the
constraints listed.

\begin{problems}
  \item Solve $\max\ z=2x_{1}+3x_{2}$ subject to $x_{1}+x_{2}\le 5$.
  \item Find all corner points of
        $\{x_{1}+3x_{2}\le 9,\ x_{1}+x_{2}\le 5,\ x_{1},x_{2}\ge 0\}$ and
        maximize $z=4x_{1}+5x_{2}$ over this region.
  \item Solve $\max\ z=2x_{1}+2x_{2}$ subject to $x_{1}+x_{2}\le 6$ and
        $x_{1}\le 5$, and describe the complete set of optimal solutions.
  \item Decide whether $\max\ z=3x_{1}+x_{2}$ subject to
        $x_{1}-2x_{2}\le 2$ has a finite optimum. If not, exhibit a
        feasible ray along which $z\to\infty$.
  \item Show that the constraint set $2x_{1}+x_{2}\le 4$,
        $x_{1}+x_{2}\ge 6$, $x_{1},x_{2}\ge 0$ is infeasible.
  \item Solve $\min\ z=3x_{1}+2x_{2}$ subject to $x_{1}+x_{2}\ge 6$,
        $x_{1}\ge 1$, $x_{2}\ge 1$, and classify the constraints at the
        optimum.
  \item A bakery sells bread ($x_{1}$, profit $3$) and cake ($x_{2}$,
        profit $5$). Flour: $x_{1}+2x_{2}\le 12$; sugar:
        $x_{1}+x_{2}\le 8$. Find the profit-maximizing plan.
  \item Solve $\max\ z=2x_{1}+x_{2}$ subject to $x_{1}+x_{2}\le 5$,
        $x_{1}\le 3$, $x_{2}\le 4$. Classify every constraint at the
        optimum as binding or slack, giving slack values.
\end{problems}

\section{Additional Solutions}

\begin{solutions}
  \item Corners of $\{x_{1}+x_{2}\le 5,\ x_{1},x_{2}\ge 0\}$: $(0,0)$,
        $(5,0)$, $(0,5)$. Evaluate $z=2x_{1}+3x_{2}$:
        \[
          z(0,0)=0,\qquad z(5,0)=10,\qquad z(0,5)=15 .
        \]
        Optimum $z^{*}=15$ at $(0,5)$; binding there:
        $x_{1}+x_{2}\le 5$ and $x_{1}\ge 0$.

  \item Label (C1) $x_{1}+3x_{2}=9$ and (C2) $x_{1}+x_{2}=5$.
        Candidates: $(0,0)$ feasible; $x_{2}=0$ with (C1): $(9,0)$,
        check (C2): $9>5$, infeasible; $x_{2}=0$ with (C2): $(5,0)$,
        check (C1): $5\le 9$, corner; $x_{1}=0$ with (C1): $(0,3)$,
        check (C2): $3\le 5$, corner; $x_{1}=0$ with (C2): $(0,5)$,
        check (C1): $15>9$, infeasible; (C1) with (C2): subtracting
        (C2) from (C1) gives $2x_{2}=4$, so $x_{2}=2$, $x_{1}=3$,
        corner. Corners: $(0,0)$, $(5,0)$, $(3,2)$, $(0,3)$. Evaluate
        $z=4x_{1}+5x_{2}$:
        \[
          z(0,0)=0,\quad z(5,0)=20,\quad z(3,2)=12+10=22,\quad
          z(0,3)=15 .
        \]
        Optimum $z^{*}=22$ at $(3,2)$; binding: (C1) ($3+6=9$) and
        (C2) ($3+2=5$).

  \item Label (C1) $x_{1}+x_{2}\le 6$, (C2) $x_{1}\le 5$. Candidates:
        $(0,0)$ feasible; (C2) with $x_{2}=0$: $(5,0)$, check (C1):
        $5\le 6$, corner; (C1) with $x_{2}=0$: $(6,0)$, violates (C2),
        infeasible; (C1) with $x_{1}=0$: $(0,6)$, check (C2):
        $0\le 5$, corner; (C1) with (C2): $(5,1)$, feasible, corner.
        Corners: $(0,0)$, $(5,0)$, $(5,1)$, $(0,6)$. Evaluate
        $z=2x_{1}+2x_{2}$:
        \[
          z(0,0)=0,\quad z(5,0)=10,\quad z(5,1)=12,\quad z(0,6)=12 .
        \]
        The maximum $z^{*}=12$ is attained at both $(5,1)$ and
        $(0,6)$: the iso-profit lines $2x_{1}+2x_{2}=z$ are parallel
        to the binding line $x_{1}+x_{2}=6$. Optimal set: the edge
        $\{(5-5t,\,1+5t):t\in[0,1]\}$, on which $x_{1}+x_{2}=6$
        throughout.

  \item The region $\{x_{1}-2x_{2}\le 2,\ x_{1},x_{2}\ge 0\}$ has
        corners $(0,0)$ and $(2,0)$ (the pair $x_{1}=0$ with
        $x_{1}-2x_{2}=2$ gives $(0,-1)$, infeasible), but it is
        unbounded: for every $t\ge 0$ the point $(0,t)$ satisfies
        $0-2t\le 2$. Along this ray $z=3(0)+t=t\to\infty$, so the LP
        is unbounded and has no finite optimum.

  \item For any point with $x_{1},x_{2}\ge 0$,
        \[
          x_{1}+x_{2}\;\le\;2x_{1}+x_{2}\;\le\;4,
        \]
        which contradicts $x_{1}+x_{2}\ge 6$. Hence no feasible point
        exists: the LP is infeasible (the half-plane on or below
        $2x_{1}+x_{2}=4$ lies entirely below the line $x_{1}+x_{2}=6$
        within the first quadrant).

  \item Label (C1) $x_{1}+x_{2}\ge 6$, (C2) $x_{1}\ge 1$, (C3)
        $x_{2}\ge 1$. Candidates: (C1) with (C2): $(1,5)$, check
        (C3): $5\ge 1$, corner; (C1) with (C3): $(5,1)$, check (C2):
        $5\ge 1$, corner; (C2) with (C3): $(1,1)$, check (C1):
        $2<6$, infeasible. The region is unbounded, but both
        coefficients of $z=3x_{1}+2x_{2}$ are positive, so moving
        deeper into the region only increases $z$; the minimum sits at
        a corner. Evaluate:
        \[
          z(1,5)=3+10=13,\qquad z(5,1)=15+2=17 .
        \]
        Minimum $z^{*}=13$ at $(1,5)$. Binding there: (C1)
        ($1+5=6$) and (C2) ($x_{1}=1$); slack: (C3) with slack value
        $5-1=4$.

  \item Model: $\max\ z=3x_{1}+5x_{2}$ subject to flour (C1)
        $x_{1}+2x_{2}\le 12$, sugar (C2) $x_{1}+x_{2}\le 8$.
        Candidates: $(0,0)$ feasible; $x_{2}=0$ with (C1): $(12,0)$,
        check (C2): $12>8$, infeasible; $x_{2}=0$ with (C2): $(8,0)$,
        check (C1): $8\le 12$, corner; $x_{1}=0$ with (C1): $(0,6)$,
        check (C2): $6\le 8$, corner; $x_{1}=0$ with (C2): $(0,8)$,
        check (C1): $16>12$, infeasible; (C1) with (C2): subtracting
        gives $x_{2}=4$, hence $x_{1}=4$, corner. Corners: $(0,0)$,
        $(8,0)$, $(4,4)$, $(0,6)$. Profits:
        \[
          z(0,0)=0,\quad z(8,0)=24,\quad z(4,4)=12+20=32,\quad
          z(0,6)=30 .
        \]
        Optimum: $4$ breads and $4$ cakes with profit $z^{*}=32$; both
        flour ($4+8=12$) and sugar ($4+4=8$) are binding.

  \item Label (C1) $x_{1}+x_{2}\le 5$, (C2) $x_{1}\le 3$, (C3)
        $x_{2}\le 4$. Candidates: $(0,0)$ feasible; (C2) with
        $x_{2}=0$: $(3,0)$, check (C1): $3\le 5$, corner; (C1) with
        $x_{2}=0$: $(5,0)$, violates (C2), infeasible; (C3) with
        $x_{1}=0$: $(0,4)$, check (C1): $4\le 5$, corner; (C1) with
        $x_{1}=0$: $(0,5)$, violates (C3), infeasible; (C1) with
        (C2): $(3,2)$, check (C3): $2\le 4$, corner; (C1) with (C3):
        $(1,4)$, check (C2): $1\le 3$, corner; (C2) with (C3):
        $(3,4)$, check (C1): $7>5$, infeasible. Corners: $(0,0)$,
        $(3,0)$, $(3,2)$, $(1,4)$, $(0,4)$. Evaluate
        $z=2x_{1}+x_{2}$:
        \[
          z(0,0)=0,\quad z(3,0)=6,\quad z(3,2)=8,\quad z(1,4)=6,\quad
          z(0,4)=4 .
        \]
        Optimum $z^{*}=8$ at $(3,2)$. Classification: (C1) binding
        ($3+2=5$, slack $0$); (C2) binding ($x_{1}=3$, slack $0$);
        (C3) slack with slack value $4-2=2$; both sign constraints
        slack.
\end{solutions}
